% Farben für jede Schicht (\providecommand: auch im Skript definiert)
\providecommand{\applicationLayerColor}{Gold}
\providecommand{\transportLayerColor}{Indigo}
\providecommand{\networkLayerColor}{RoyalBlue}
\providecommand{\dataLinkLayerColor}{Gray}

% command to make sure colorboxes don't create whitespace between code lines
\fboxsep=.6pt

% default settings for listings
\lstset{style=python}

\title[Netzwerke]{Netzwerke}

\subtitle{Anwendungsschicht (Skript, Abschnitte 1.6 bis 1.8)}

\author[Blum, Cyril] % (optional)
{Cyril Blum}

\institute[KLW] % (optional)
{
    Fachschaft Informatik\\
    Kantonsschule im Lee
}

\date[\today] % (optional)

\logo{\includegraphics[width=2cm]{Logos/FreeFlower.pdf}}

\titlegraphic{\vspace{-.7cm}\includegraphics[width=.85\textwidth]{"Grundlagen_Info/06_Datenbanken/Figures/title-emojis"}}

\frame{\titlepage}

\section{Anwendungsschicht}
\setbeamercolor{frametitle}{bg=\applicationLayerColor!20}

\begin{frame}{Netzwerk-Spiel: \textbf{Anwendungsschicht}}
    Endgeräte verfassen Nachrichten:
    \begin{table}[H]
        \centering
        \begin{tabular}{p{2cm}l}
            \rowcolor{\applicationLayerColor!20}Nachricht\newline\tiny (max. 10 Zeichen): & \texttt{Hi Bob, wi} \\
        \end{tabular}
        \begin{tabular}{p{2cm}l}
            \rowcolor{\applicationLayerColor!20}Nachricht\newline\tiny (max. 10 Zeichen): & \texttt{e geht's?} \\
        \end{tabular}
    \end{table}
\end{frame}


\begin{frame}[fragile]{\faLaptop{} Client vs Server vs Host}{Anwendungsschicht}
    Definitionen
    \begin{itemize}[<+->]
        \item \textbf{Client}
              \begin{itemize}
                  \item \enquote{Kunde} (von en. \textit{client})
                  \item Computer(-Programm), das etwas will
              \end{itemize}
        \item \textbf{Server}
              \begin{itemize}
                  \item \enquote{Diener} (von en. \textit{to serve} = bedienen, $\neq$ \enquote{surf})
                  \item Computer(-Programm), das etwas hat, das andere Computer(-Programme) brauchen
                  \item Kann Verbindungsanfragen von anderen Computer(-Programmen) akzeptieren / ablehnen
              \end{itemize}
        \item \textbf{Host}
              \begin{itemize}
                  \item = Computer in einem Netzwerk, allgemein für \enquote{server} oder \enquote{client}
                  \item \textit{Host name} = Name eines Computers
              \end{itemize}

    \end{itemize}
\end{frame}

\begin{frame}[fragile]{\faSafari{} Browser}{Anwendungsschicht}
    \begin{itemize}
        \item Browser = Programm, kann eine Datei (meist \texttt{.html})...
              \begin{itemize}
                  \item ...über das Internet \textit{anfordern}
                  \item ...nach dem Empfang \textit{darstellen}
              \end{itemize}
        \item Das Senden / Empfangen von Dateien übernimmt das Betriebssystem (z.B. Windows, MacOS)
        \item Beispiele:
              \begin{itemize}
                  \item \faChrome{} Chrome
                  \item \faFirefox{} Firefox
                  \item \faSafari{} Safari
                  \item \faEdge{} Edge
                  \item  etc.
              \end{itemize}
    \end{itemize}
\end{frame}

\begin{frame}[fragile]{\acrfull{DNS}}
    \begin{itemize}
        \item Wenn Sie auf eine Webseite wie \url{https://sbb.ch} gehen, woher weiss der Client, von welchem Server er die Informationen anfragen soll?
        \item \gls{DNS}: Domain-Name $\rightleftarrows$ \acrshort{IP}-Adresse!
    \end{itemize}
\end{frame}

\begin{frame}[fragile]{\gls{DNS}}
    \begin{figure}[H]
        \resizebox{\linewidth}{!}{
            \begin{tikzpicture}[
                    base/.style={rectangle, draw, minimum width=2cm, minimum height=1cm, text=white, rounded corners, text width=2.3cm, align=center,inner sep=8pt},
                    server/.style={base, fill=\networkLayerColor},
                    client/.style={base, fill=ForestGreen},
                    arrow/.style={-Stealth, thick, rounded corners,text width=4cm,align=center}
                ]

                % Nodes
                \node[client] (client) {Client\\\texttt{192.168.0.10}};
                \node[server, right=4.2cm of client] (dns) {\gls{DNS}-Server\\\texttt{192.168.2.1}};
                \node[server, right=4.2cm of dns] (web) {Webserver\\\texttt{192.168.0.11}};
                \node[text width=3cm, right=0cm of dns] (dnsTable) {\tiny \begin{tabular}{p{1.1cm}p{1.2cm}}
                        \textbf{Domain\newline Name} & \textbf{\acrshort{IP}-\newline Adresse} \\
                        \midrule
                        \texttt{beispiel.ch}         & 192.168.0.11                  \\
                        \midrule
                        \texttt{schule.ch}           & 192.168.0.12                  \\
                        \midrule
                        \texttt{test.org}            & 192.168.0.13                  \\
                        \midrule
                        \texttt{insta.com}           & 192.168.0.14                  \\
                        \midrule
                        \texttt{...}                 & ...                           \\
                    \end{tabular}};

                % draw box around DNS node and DNS table
                \node[fit=(dns) (dnsTable), inner sep=4pt, \networkLayerColor, dotted, draw, line width=1pt,rounded corners](dnsArea) {};

                % Arrows
                \draw[arrow] ($(client.east)+(0,0.15)$) -- ($(dns.west)+(0,0.15)$) node[midway, above] {\textbf{1.} \gls{DNS}-Anfrage\\\footnotesize\enquote{Wer ist \texttt{beispiel.ch}?}};
                \draw[arrow] ($(dns.west)-(0,0.15)$) -- ($(client.east)-(0,0.15)$) node[midway, below] {\textbf{2.} \gls{DNS}-Antwort\\\footnotesize\enquote{\texttt{192.168.0.11}}};
                \draw[arrow] (client) -- ++(0,2cm) -- node[above] {\textbf{3.} \gls{HTTP}-Anfrage\\\footnotesize (Webseite wird angefragt)} ++($(web)$) -- (web);
                \draw[arrow] (web) -- ++(0,-2cm) -- node[below,midway] {\textbf{4.} \gls{HTTP}-Antwort\\\footnotesize (Webseite wird geliefert)} ++($(client)-(web)$) -- (client);

            \end{tikzpicture}}
        \centering
        \caption{Schematischer Ablauf einer \gls{DNS}-Anfrage}
        \label{fig:dns}
    \end{figure}
\end{frame}

\begin{frame}[fragile]{\faSitemap{} \acrshort{URL} (Uniform Resource Locator)}{Präsentationsschicht: Webseiten}

    \acrshort{URL} = Web-Adresse

    \[
        \underbrace{\text{http://}}_{\text{Protokoll}}
        \underbrace{\text{www.}}_{\text{Subdomain}}
        \underbrace{\text{beispiel.}}_{\text{Domain}}
        \underbrace{\text{ch}}_{\text{\acrshort{TLD}\tikzmark{tld}}}
        \underbrace{\text{/dokumente/}}_{\text{Ordner}}
        \underbrace{\text{reglemente.html}}_{\text{Dateiname}}
    \]
    \begin{tikzpicture}[overlay, remember picture]
        \node[below=of tld] (tldFull) {Top Level Domain};
        \draw[red,-stealth] (tld.north) to[bend right] (tldFull);
    \end{tikzpicture}

    \begin{itemize}[<+->]
        \item Erster Teil bezeichnet einen \textit{Server}
              \begin{itemize}
                  \item z.B. \texttt{www.beispiel.ch} $\rightarrow$ 58.55.110.9
              \end{itemize}
        \item Alles nach der \acrshort{TLD} ist ein normaler Ordnerpfad auf dem Computer
    \end{itemize}

\end{frame}

\begin{frame}[fragile]{\faServer{} \acrshort{HTTP}(S)-Protokoll}
    \begin{figure}[H]
        \centering
        \resizebox{\linewidth}{!}{
            \begin{tikzpicture}[
                    font=\sffamily,
                    every node/.style={align=center},
                    entity/.style={rectangle, draw, fill=\networkLayerColor, text=white, minimum height=2.5cm, minimum width=3cm, rounded corners, text width=2.3cm, align=center, node distance=1.5cm and 3.5cm},
                    process/.style={-Stealth, thick, draw},
                    arrowtext/.style={text width=2.5cm, align=center, text=black, font=\footnotesize}
                ]

                % Nodes for user
                \node[entity] (userStart) {Eingabe der \acrshort{URL} am Bildschirm};
                \node[entity, below=of userStart] (userEnd) {Webseite am Bildschirm};
                % Draw a 2D rectangle around the users
                \node[fit=(userStart) (userEnd), inner sep=4pt, red, text=black, dotted, draw, line width=1pt,rounded corners](user) {User \\ \faUser};

                % nodes for
                \node[entity, right=of userStart] (browserStart) {Übersetzt die \acrshort{URL} in einen \acrshort{HTTP}-Request};
                \node[entity, right=of userEnd] (browserEnd) {Sammelt Daten, setzt Website zusammen (rendern)};
                % Draw a 2D rectangle around the users
                \node[fit=(browserStart) (browserEnd), inner sep=4pt, red, text=black, dotted, draw, line width=1pt,rounded corners](browser) {Browser\\ \faLaptop};

                % nodes for server
                \node[entity, right=of browserStart] (serverStart) {Interpretiert den Request und sucht die Dateien};
                \node[entity, right=of browserEnd] (serverEnd) {Sendet den Status der Suche und die Daten zur Webseite};
                % Draw a 2D rectangle around the users
                \node[fit=(serverStart) (serverEnd), inner sep=4pt, red, text=black, dotted, draw, line width=1pt,rounded corners](server) {Web Server\\ \faServer};

                % Connections
                \draw[process, Crimson] (userStart.east) -- node[arrowtext, above] {1. \acrshort{URL} \\ \texttt{http://\\www.beispiel.ch/\\index.html}} (browserStart.west);
                \draw[process, Crimson] (browserStart.east) -- node[arrowtext, above] {2. \acrshort{HTTP} Request \\ \texttt{GET /index.html}} (serverStart.west);
                \draw[process, ForestGreen] (serverEnd.west) -- node[arrowtext, below] {3. \acrshort{HTTP} Response \\ Statuscode + Daten} (browserEnd.east);
                \draw[process, ForestGreen] (browserEnd.west) -- node[arrowtext, below] {4. Webpage \\ Darstellung der Webseite \texttt{index.html}} (userEnd.east);

            \end{tikzpicture}}
    \end{figure}
\end{frame}



\begin{frame}[fragile]{\faLaptop{} Webseiten}{Präsentationsschicht}
    Beispiel: \acrshort{HTML} (\textit{Hypertext Markup Language})
    \begin{itemize}[<+->]
        \item Demo (\textit{Inspect Element} oder \textit{View Page Source})
        \item Speichert die Struktur einer Webseite
        \item Kann weitere Dateien enthalten, beispielsweise:
              \begin{itemize}
                  \item CSS (Cascading Style Sheet), für die Darstellung (Farbe, Grösse etc.)
                  \item JavaScript, jQuery etc. für die Interaktivität (Ein- und Ausblenden von Objekten etc.)
              \end{itemize}
        \item Wie man eine Webseite baut ist hier nicht das Hauptthema, wird aber in der Interessenwoche vermittelt
    \end{itemize}
\end{frame}

\begin{frame}[fragile]{Web- und \acrshort{DNS}-Server}{Aufträge (Skript)}
    \textbf{Webserver \textit{ohne} \gls{DNS}}
    \begin{itemize}
        \item \aufgabebox{Filius (Web-Server): 1.35}
        \item \challengebox{1.36 (\acrshort{HTML} besser verstehen)}
        \item \challengebox{Filius (\enquote{Echter} Web-Server): 1.37}
    \end{itemize}
    \textbf{\gls{DNS}}
    \begin{itemize}
        \item \aufgabebox{1.38 - 1.40}
    \end{itemize}
\end{frame}

\section{Zusammenfassung}
\setbeamercolor{frametitle}{bg=white}

\begin{frame}{Ablauf eines Webaufrufs}{Alle Schichten arbeiten zusammen}
    \begin{center}
        \adjustbox{max width=\textwidth, max totalheight=.8\textheight}{% Ablauf eines Webaufrufs als Sequenzdiagramm (Laptop, Router, DNS-Server, Webserver).
% Farben der Pfeile = Schicht des jeweiligen Schritts (Farben aus Skript.tex bzw. den Folien).
% Adressen aus den Dokumentationsbereichen (RFC 5737), damit keine echten Hosts gemeint sind.
% \glsentryshort statt \acrshort: keine farbigen Glossar-Links in der Grafik.
\begin{tikzpicture}[>=Stealth, font=\footnotesize,
        actor/.style={draw, rounded corners, fill=ForestGreen!20, align=center, minimum width=2.6cm, minimum height=1cm},
        stepno/.style={circle, draw=#1, text=#1, fill=white, thick, inner sep=1.5pt, font=\scriptsize\bfseries},
        stepno/.default=black,
        msg/.style={->, thick},
        note/.style={draw=#1, text=#1, fill=white, rounded corners, align=center, font=\scriptsize, inner sep=3pt},
        note/.default=black
    ]
    \def\appCol{\applicationLayerColor!80!black}

    % Teilnehmende (x-Positionen)
    \node[actor] (pc)  at (0,0)    {\faLaptop{} Laptop (Browser)\\\texttt{192.168.1.23}};
    \node[actor] (rt)  at (3.8,0)  {Router / Gateway\\\texttt{192.168.1.1}};
    \node[actor] (dns) at (7.6,0)  {\glsentryshort{DNS}-Server\\\texttt{203.0.113.53}};
    \node[actor] (web) at (11.4,0) {Webserver\\\texttt{198.51.100.7}};

    % Zeitachsen
    \foreach \n in {pc, rt, dns, web} {
            \draw[gray, dashed] (\n.south) -- ++(0,-11.2);
        }

    % 1: ARP (Netzzugangsschicht)
    \draw[msg, \dataLinkLayerColor] (0,-1.4) -- node[above] {Wer hat \texttt{192.168.1.1}?} (3.8,-1.4);
    \draw[msg, \dataLinkLayerColor] (3.8,-2.1) -- node[above] {Ich, \glsentryshort{MAC} \texttt{a4:2b:\dots}} (0,-2.1);
    \node[stepno=\dataLinkLayerColor] at (-1.6,-1.75) {1};

    % 2: DNS (Anwendungsschicht, über UDP)
    \draw[msg, \appCol] (0,-3.0) -- node[above, pos=.3] {\glsentryshort{IP} von \texttt{www.example.org}?} (7.6,-3.0);
    \draw[msg, \appCol] (7.6,-3.7) -- node[above, pos=.7] {\texttt{198.51.100.7}} (0,-3.7);
    \node[stepno=\appCol] at (-1.6,-3.35) {2};

    % 3: NAT und Routing (Vermittlungsschicht)
    \node[note=\networkLayerColor] at (3.8,-4.6) {\glsentryshort{NAT}: Absender wird\\\texttt{203.0.113.45} (öffentlich)};
    \node[note=\networkLayerColor] at (7.6,-4.6) {Routing über viele\\Router im Internet};
    \node[stepno=\networkLayerColor] at (-1.6,-4.6) {3};

    % 4: TCP-Verbindungsaufbau (Transportschicht)
    \draw[msg, \transportLayerColor] (0,-5.6) -- node[above, pos=.2] {\texttt{SYN} (an Port 443)} (11.4,-5.6);
    \draw[msg, \transportLayerColor] (11.4,-6.2) -- node[above, pos=.8] {\texttt{SYN-ACK}} (0,-6.2);
    \draw[msg, \transportLayerColor] (0,-6.8) -- node[above, pos=.2] {\texttt{ACK}} (11.4,-6.8);
    \node[stepno=\transportLayerColor] at (-1.6,-6.2) {4};

    % 5: TLS (Anwendungsschicht)
    \draw[<->, thick, \appCol] (0,-7.7) -- node[above, pos=.35] {\glsentryshort{TLS}: Zertifikat prüfen, Schlüssel vereinbaren} (11.4,-7.7);
    \node[stepno=\appCol] at (-1.6,-7.7) {5};

    % 6: HTTP (Anwendungsschicht)
    \draw[msg, \appCol] (0,-8.7) -- node[above, pos=.25] {\texttt{GET /index.html}} (11.4,-8.7);
    \draw[msg, \appCol] (11.4,-9.4) -- node[above, pos=.75] {\texttt{200 OK} + \glsentryshort{HTML}-Datei} (0,-9.4);
    \node[stepno=\appCol] at (-1.6,-9.05) {6};

    % 7: Darstellung und weitere Anfragen
    \node[note, anchor=west] at (-0.9,-10.6) {Seite anzeigen, Bilder \& Stylesheets nachladen};
    \node[stepno] at (-1.6,-10.6) {7};
\end{tikzpicture}
}
    \end{center}
\end{frame}

\begin{frame}{Ablauf eines Webaufrufs}{Schritt für Schritt}
    \begin{enumerate}[<+->]\small
        \item \textcolor{\dataLinkLayerColor}{\textbf{Netzzugang}}: \acrshort{MAC}-Adresse des Routers mit \acrshort{ARP} erfragen
        \item \textcolor{\applicationLayerColor!80!black}{\textbf{Anwendung}}: \acrshort{IP}-Adresse des Webservers mit \acrshort{DNS} erfragen
        \item \textcolor{\networkLayerColor}{\textbf{Vermittlung}}: \acrshort{NAT} im Heimrouter, Routing durchs Internet
        \item \textcolor{\transportLayerColor}{\textbf{Transport}}: \acrshort{TCP}-Verbindung zu Port 443 (3-way handshake)
        \item \textcolor{\applicationLayerColor!80!black}{\textbf{Anwendung}}: \acrshort{TLS} -- Zertifikat prüfen, Schlüssel vereinbaren
        \item \textcolor{\applicationLayerColor!80!black}{\textbf{Anwendung}}: \acrshort{HTTP}-Anfrage \texttt{GET}, Antwort \texttt{200 OK} mit \acrshort{HTML}
        \item Seite darstellen, Bilder \& Stylesheets nachladen
    \end{enumerate}
\end{frame}

\begin{frame}{Auftrag (Skript)}
    \begin{itemize}
        \item \aufgabebox{1.41 (Ablauf eines Webaufrufs)}
        \item \challengebox{1.42 (Webaufruf in Wireshark)}
        \item \challengebox{1.45 - 1.46 (Filius-Troubleshooting, Vertiefung Schichtenmodell)}
    \end{itemize}
\end{frame}
